\documentclass[10pt]{article}
\usepackage[utf8]{inputenc}
\usepackage[T1]{fontenc}
\usepackage[spanish,es-tabla]{babel}
\usepackage[margin=2cm]{geometry}
\usepackage{booktabs,float,url,hyperref}
\title{\large Anexo independiente: oclusión en 582 casos aceptados de fauna V6}
\author{Equipo de la Pregunta 3 -- Universidad Nacional de Ingeniería}
\date{8 de octubre de 2026}
\begin{document}
\maketitle

\section{Cohorte y decisiones preservadas}
La muestra original contiene 600 objetos de los 17\,620 elegibles de trans\_test.
La revisión aceptada conserva 582: 340 con oclusión baja, 69 media, 48 alta y
125 ambiguos. Se excluyeron los 18 sin nivel por decisión expresa del responsable.
El caso 121 (anotación 23098) mantiene el nivel medium según Miguel y la nota vacía.
Las notas se aceptaron como opcionales: 398 permanecen ausentes. No se inventan
justificaciones ni niveles; los retornos humanos y la discrepancia original se preservan.
Los casos excluidos son 253, 302, 385--393, 409--414 y 439.

Baja corresponde a más de 75\,\% del cuerpo visible; media, 25--75\,\%; alta,
menos de 25\,\%. La visibilidad incierta o el truncamiento se conserva como ambiguous.
Cada nivel evaluado supera el mínimo protocolario de 30 objetos. Los ambiguos
no constituyen un cuarto grupo de AP.

\section{Resultados y método}
\begin{table}[H]
\centering\small
\caption{Oclusión en la submuestra aceptada de trans\_test. AP no ponderada sobre todos sus fotogramas; IC95 por bootstrap de secuencias, 200 réplicas por nivel. Los 125 ambiguos no son un cuarto grupo de AP.}
\begin{tabular}{lrrr}
\toprule
Nivel & Cajas & mAP50 [IC95] & mAP50:95 [IC95] \\
\midrule
Baja & 340 & 0.1582 [0.1457, 0.2117] & 0.0875 [0.0809, 0.1217] \\
Media & 69 & 0.0582 [0.0389, 0.1035] & 0.0271 [0.0184, 0.0523] \\
Alta & 48 & 0.0563 [0.0329, 0.1164] & 0.0244 [0.0160, 0.0567] \\
\bottomrule
\end{tabular}
\end{table}

Se reutilizan el checkpoint y las predicciones congeladas V6, sin entrenamiento
ni inferencia nuevos. Se mantienen los 18\,647 fotogramas de trans\_test, incluidos
1\,778 vacíos. Las cajas no revisadas, excluidas, ambiguas o de otro nivel son
regiones crowd ignoradas según las reglas COCO congeladas. Las coincidencias de
la misma clase pueden ignorarse; las detecciones no absorbidas conservan la
posibilidad de ser falsos positivos, también en los fotogramas vacíos.
No se calculó AP únicamente sobre las fotografías revisadas.

La AP no aplica pesos de muestreo. Los intervalos percentiles del 95\,\% usan
200 remuestreos de las 7\,438 secuencias por nivel, semilla 42 y checkpoint fijo.
Las 200 réplicas de ambas métricas quedaron definidas en cada nivel.
Los intervalos son condicionales a las secuencias y cámaras observadas; no
incluyen desacuerdo de anotación, notas faltantes, semillas de entrenamiento
ni rendimiento en cámaras futuras.

\medskip
\section{Interpretación y alcance descriptivo}
La mAP50:95 baja 6.04 puntos porcentuales de baja a media y 6.31 de baja a alta.
Media y alta difieren 0.28 puntos, con intervalos ampliamente superpuestos; no se
calculó un intervalo de la diferencia. Los grupos tienen 13, 10 y 11 categorías
con soporte: no se aisló el efecto causal de oclusión. La disminución de positivos
evaluables junto con la conservación de fotogramas y falsos positivos cambia
la tarea de evaluación. Esta mAP no es directamente comparable con la mAP global V6.

Las proporciones entre aceptados son 58.42\,\% baja, 11.86\,\% media,
8.25\,\% alta y 21.48\,\% ambigua. Las razones de pesos originales conservados
son 59.75\,\%, 13.02\,\%, 7.55\,\% y 19.68\,\%, respectivamente.
Los excluidos representan 2.38\,\% del peso de diseño original, pero son falta
de respuesta humana, no una nueva muestra aleatoria. No se reutilizan los
intervalos poblacionales de proporciones ni la garantía de planificación
de 600 casos; no se afirma representatividad poblacional posterior a las exclusiones.

\newpage
\section{Detalle por especie y trazabilidad}
Las tablas siguientes conservan AP y soporte de los resultados aceptados.
ND significa ausencia de referencia evaluable, no ausencia biológica ni AP cero.
El CSV numérico y ambos notebooks acompañan el informe. Los valores sin redondeo
y los IC están en \texttt{data/results.json}.
\begin{table}[H]
\centering\small
\caption{Soporte y AP50:95 por especie y nivel de oclusión.}
\begin{tabular}{lr@{\hspace{1em}}rr@{\hspace{1em}}rr@{\hspace{1em}}r}
\toprule
 & \multicolumn{2}{c}{Baja} & \multicolumn{2}{c}{Media} & \multicolumn{2}{c}{Alta} \\
Especie & Cajas & AP50:95 & Cajas & AP50:95 & Cajas & AP50:95 \\
\midrule
opossum & 59 & 0.3104 & 20 & 0.1204 & 8 & 0.1551 \\
raccoon & 66 & 0.2459 & 10 & 0.1276 & 7 & 0.0258 \\
squirrel & 19 & 0.0318 & 7 & 0.0015 & 3 & 0.0003 \\
bobcat & 42 & 0.1495 & 3 & 0.0055 & 5 & 0.0190 \\
skunk & 15 & 0.0336 & 0 & ND & 5 & 0.0154 \\
dog & 27 & 0.1633 & 4 & 0.0078 & 3 & 0.0475 \\
coyote & 31 & 0.0599 & 9 & 0.0032 & 5 & 0.0019 \\
rabbit & 24 & 0.0682 & 2 & 0.0019 & 3 & 0.0031 \\
bird & 21 & 0.0384 & 3 & 0.0023 & 1 & 0.0000 \\
cat & 24 & 0.0250 & 10 & 0.0012 & 6 & 0.0000 \\
badger & 5 & 0.0048 & 1 & 0.0000 & 0 & ND \\
deer & 0 & ND & 0 & ND & 0 & ND \\
fox & 1 & 0.0000 & 0 & ND & 0 & ND \\
rodent & 6 & 0.0069 & 0 & ND & 2 & 0.0000 \\
\bottomrule
\end{tabular}
\end{table}
\begin{table}[H]
\centering\small
\caption{AP50 por especie y nivel de oclusión; soportes en la tabla anterior.}
\begin{tabular}{lrrr}
\toprule
Especie & Baja & Media & Alta \\
\midrule
opossum & 0.7330 & 0.3246 & 0.4354 \\
raccoon & 0.4355 & 0.2133 & 0.0430 \\
squirrel & 0.0509 & 0.0043 & 0.0022 \\
bobcat & 0.1972 & 0.0079 & 0.0241 \\
skunk & 0.0437 & ND & 0.0220 \\
dog & 0.2772 & 0.0186 & 0.0849 \\
coyote & 0.0808 & 0.0046 & 0.0025 \\
rabbit & 0.1140 & 0.0033 & 0.0052 \\
bird & 0.0756 & 0.0031 & 0.0000 \\
cat & 0.0319 & 0.0020 & 0.0000 \\
badger & 0.0073 & 0.0000 & ND \\
deer & ND & ND & ND \\
fox & 0.0000 & ND & ND \\
rodent & 0.0092 & ND & 0.0000 \\
\bottomrule
\end{tabular}
\end{table}


La firma del análisis aceptado es
\begin{quote}\small\nolinkurl{6b2983c33abccbc18e293cc3a10ca30a710669c2ba8055a76bb462b3348a13a9}\end{quote}
El SHA-256 del JSON de resultados es
\begin{quote}\small\nolinkurl{b5ede4e0f8d46aef5b579f64c3b7f7c476b4aeb93eafaeffaafcdd4e0da124f9}\end{quote}
\texttt{integridad\_datos.json} permite verificar las copias byte a byte. Los
originales V6 permanecen intactos. Este anexo no altera sus resultados globales,
por tamaño, Grad-CAM ni monitoreo. No se atribuyen aprobación del docente ni
envío a la plataforma del curso.
\end{document}
